% Generated draft: requires mathematical and source review.
\documentclass[11pt]{article}
\usepackage[margin=25mm]{geometry}
\usepackage{fontspec}
\setmainfont{Latin Modern Roman}
\usepackage{amsmath,amssymb,mathtools,unicode-math}
\setmathfont{Latin Modern Math}
\usepackage{enumitem,needspace,xurl,hyperref}
\setlength{\parindent}{0pt}
\setlength{\parskip}{5pt}
\setlength{\emergencystretch}{3em}
\newcommand{\N}{\mathbb{N}}
\newcommand{\Z}{\mathbb{Z}}
\newcommand{\Q}{\mathbb{Q}}
\newcommand{\R}{\mathbb{R}}
\newcommand{\C}{\mathbb{C}}
\newcommand{\F}{\mathbb{F}}
\newcommand{\A}{\mathbb{A}}
\newcommand{\PP}{\mathbb{P}}
\newcommand{\E}{\mathbb{E}}
\newcommand{\eps}{\varepsilon}
\newcommand{\dd}{\,\mathrm{d}}
\begin{document}
\section*{UnsolvedMath ID: 30004062}
\textbf{Source problem code:} catalogue entry
\subsection{Definitions and mathematical statement}
The supplied record poses a research question about reconstructing functions with discontinuities from discrete tomographic measurements. For schematic notation only, write
\[
\widehat f=\mathcal R\bigl(\mathcal D(\mathcal T f)\bigr),
\]
where \(f\) is the unknown function, \(\mathcal T\) denotes a tomographic measurement operator, \(\mathcal D\) denotes discrete sampling of its data, and \(\mathcal R\) denotes a reconstruction procedure. These symbols are placeholders: the source record does not specify their domains, formulas, admissibility conditions, or sampling parameters.

\begin{enumerate}
\item Describe the behavior of \(\widehat f\) near singularities of \(f\).
\item Determine how singularities of \(f\) remote from a reconstruction location affect the reconstruction at that location.
\item Explain under what conditions, and in what mathematical sense, functions with discontinuities can be reconstructed accurately from discrete tomographic data.
\end{enumerate}

The supplied formulation does not specify the spatial dimension or domain, the class of functions or singularities, the measurement geometry, the reconstruction algorithm, the treatment of noise, or boundary conditions. It also does not define proximity, remoteness, accuracy, an error norm, or a limiting sampling regime. Consequently, these questions cannot be sharpened into a uniquely determined quantified theorem or conjecture from the supplied evidence. In particular, the question about accurate reconstruction does not assert accurate reconstruction for every discontinuous function or every discrete acquisition scheme.
\subsection{Short English statement}
The problem asks how discrete tomographic reconstructions behave near singularities, how distant singularities affect them, and what explains accurate reconstruction of discontinuous functions. The supplied record leaves the mathematical model and the meaning of accuracy unspecified, so a precise theorem requires additional hypotheses.
\subsection{Sources}
\begin{itemize}
\item Supplied local UnsolvedMath registry record, permanent ID \(30004062\), source problem code \texttt{catalogue entry}, title: Tomographic Reconstruction near Discontinuities. The collection is Oberwolfach Reports Open Problems. No individual proposer is identified; the supplied proposed year is 2019.
\item Supplied report citation: \emph{Tomographic Inverse Problems: Theory and Applications} (2020), Oberwolfach Reports, pp.\ 209--303, DOI 10.4171/owr/2019/4. \url{https://doi.org/10.4171/owr/2019/4}. The supplied background describes Oberwolfach Reports as documenting mathematical workshops and their open-problem discussions. No exact problem location within the cited page range is supplied.
\item Supplied literature reference: Katsevich, nonsmooth edges in discrete Radon data. \url{https://doi.org/10.1137/21M1466712}.
\item Supplied literature reference: Katsevich, three-dimensional jump discontinuities. \url{https://arxiv.org/abs/1903.08216}.
\item Supplied literature reference: Katsevich, tomographic-type reconstruction. \url{https://arxiv.org/abs/2001.05774}.
\item Literature assessment dated exactly \texttt{2026-08-22}: the supplied curation labels the problem partially solved and reports precise asymptotic descriptions near jumps and results concerning remote singularities for discrete Radon and related tomographic reconstructions in several two- and three-dimensional models, while describing the general discontinuous discrete-tomography problem as unresolved. This is supplied dated curation derived from the OWR record, not an independently verified literature assessment; its status and cited claims require independent verification.
\end{itemize}
\end{document}
